\documentclass[12pt, a4paper]{article}

\usepackage[utf8]{inputenc}
\usepackage[T1]{fontenc}
\usepackage[french]{babel}
\usepackage{amsmath, amssymb}
\usepackage{geometry}
\usepackage{hyperref}

\geometry{margin=2.5cm}

\hypersetup{colorlinks=true, linkcolor=blue, urlcolor=blue}

\title{
    \textbf{Simulation de trou noir de Schwarzschild} \\
    \large Physique relativiste et algorithme GPU
}
\author{Ascari Yannick}
\date{\today}

\begin{document}
\maketitle
\tableofcontents

\section{Introduction}

Ce document décrit les fondements physiques et algorithmiques d'une simulation
temps-réel d'un trou noir de Schwarzschild, implémentée en Rust avec le pipeline
GPU wgpu/WGSL.

L'objectif est de simuler le trajet de photons dans l'espace-temps courbé autour
d'un trou noir, en intégrant numériquement les \emph{géodésiques nulles} de la
métrique de Schwarzschild. Le résultat est rendu pixel par pixel sur le GPU, chaque
thread calculant une géodésique indépendante.

\section{La métrique de Schwarzschild}

\subsection{Définition}

La métrique de Schwarzschild décrit la géométrie de l'espace-temps autour d'une masse
sphérique statique $M$. En coordonnées sphériques $(t, r, \theta, \varphi)$, l'élément de ligne s'écrit :

\begin{equation}
    ds^2 = -\left(1 - \frac{r_s}{r}\right)c^2\,dt^2
         + \frac{dr^2}{1 - r_s/r}
         + r^2\,d\theta^2
         + r^2\sin^2\theta\,d\varphi^2
\end{equation}

Où $r_s = \frac{2GM}{c^2}$ est le \textbf{rayon de Schwarzschild}.
Dans la suite, on pose $c = G = 1$ et $r_s = 1$ (unités naturelles).

\subsection{Interprétation physique}

Le facteur $\left(1 - r_s/r\right)$ mesure la courbure locale de l'espace-temps:

\begin{itemize}
    \item $r \gg r_s$: la métrique tend vers Minkowski (espace plat).
    \item $r = r_s$: l'horizon des événements \--- aucun signal ne peut s'échapper.
    \item $r < r_s$: les rôles de $t$ et $r$ s'échangent (singularité de coordonnées). 
\end{itemize}


\section{Géodésiques nulles}

\subsection{Quantités conservées}

Un photon suit une \textbf{géodésique nulle}: $ds^2 = 0$.
La symétrie sphérique de la métrique implique, via le théorème de Noether,
deux quantités conservées le long de la géodésique:

\begin{align}
    E &= \left(1 - \frac{r_s}{r}\right)\frac{dt}{d\lambda}
    \quad \text{(énergie)} \\
    L &= r^2\frac{d\varphi}{d\lambda}
    \quad \text{(moment angulaire)}
\end{align}

Où $\lambda$ est le paramètre affine le long de la géodésique.

\subsection{Paramètre d'impact}

Le \textbf{paramètre d'impact} $b = L/E$ est la distance perpendiculaire
entre le trou noir et la trajectoire asymptotique du rayon:

\begin{equation}
    b = |\vec{r}_0 \times \hat{d}|
\end{equation}

Où $\vec{r}_0$ est la position d'origine et $\hat{d}$ la direction unitaire du rayon.

\subsection{Équation des orbites}

En posant $u = 1/r$ et en paramètrant par $\varphi$, la géodésique nulle
se réduit à:

\begin{equation}
    \boxed{\frac{d^2u}{d\varphi^2} + u = \frac{3}{2}\,r_s\,u^2}
    \label{eq:geodesique}
\end{equation}

Le terme $(-u)$ correspond à la géométrie plane \--- sans relativité, il donne
une trajectoire rectiligne. Le terme $\frac{3}{2}r_s u^2$ est la
\textbf{correction relativiste} qui courbe la lumière.

\subsection{Sphère de photons}

L'équation~\eqref{eq:geodesique} admet une orbite circulaire instable à
$r = \frac{3}{2}r_s$, la \textbf{sphère de photons}, correspondant au
paramètre d'impact critique:

\begin{equation}
    b_{\text{crit}} = \frac{3\sqrt{3}}{2}\,r_s \approx 2{,}598\,r_s
\end{equation}

Les rayons avec $b < b_{\text{crit}}$ sont absorbés. Ceux avec
$b \approx b_{\text{crit}}$ orbitent plusieurs fois, formant les
\textbf{anneaux de photons} visibles dans le rendu.

\section{Intégration numérique: méthode RK4}

\subsection{Système différentiel}

L'équation~\eqref{eq:geodesique} est du second ordre. On la convertit en
système du premier ordre en posant $v = du/d\varphi$:

\begin{equation}
    \begin{cases}
        \dfrac{du}{d\varphi} = v \\[8pt]
        \dfrac{dv}{d\varphi} = -u + \dfrac{3}{2}\,r_s\,u^2
    \end{cases}
\end{equation}

\subsection{Méthode de Runge-Kutta d'ordre 4}

On intégre ce système avec un pas $\Delta\varphi = 0{,}03\,\text{rad}$.
Pour l'état $(u_n, v_n)$ à l'étape $n$, on pose $f(u) = -u + \frac{3}{2}r_s u^2$:

\begin{align}
    k_1^u &= v_n
    & k_1^v &= f(u_n) \\
    k_2^u &= v_n + \tfrac{\Delta\varphi}{2}\,k_1^v
    & k_2^v &= f\!\left(u_n + \tfrac{\Delta\varphi}{2}\,k_1^u\right) \\
    k_3^u &= v_n + \tfrac{\Delta\varphi}{2}\,k_2^v
    & k_3^v &= f\!\left(u_n + \tfrac{\Delta\varphi}{2}\,k_2^u\right) \\
    k_4^u &= v_n + \Delta\varphi\,k_3^v
    & k_4^v &= f\!\left(u_n + \Delta\varphi\,k_3^u\right)
\end{align}

La mise à jour de l'état:

\begin{equation}
    \begin{pmatrix} u_{n+1} \\ v_{n+1} \end{pmatrix}
    = 
    \begin{pmatrix} u_n \\ v_n \end{pmatrix}
    + \frac{\Delta\varphi}{6}
    \begin{pmatrix}
        k_1^u + 2k_2^u + 2k_3^u + k_4^u \\
        k_1^v + 2k_2^v + 2k_3^v + k_4^v
    \end{pmatrix}
\end{equation}

\subsection{Conditions initiales}

Pour un rayon d'origine $\vec{r}_0$ et de direction unitaire $\hat{d}$:

\begin{align}
    u_0 &= \frac{1}{|\vec{r}_0|} \\
    \left.\frac{du}{d\varphi}\right|_0 &= -\frac{\hat{r}_0 \cdot \hat{d}}{b}
\end{align}

La position 3D à l'angle $\varphi$ est reconstruite dans la base orthonormée
du plan orbital $(\vec{e}_r, \vec{e}_\varphi)$:

\begin{equation}
    \vec{r}(\varphi) = \frac{1}{u(\varphi)}
    \left[\cos\varphi\;\vec{e}_r + \sin\varphi\;\vec{e}_\varphi\right]
\end{equation}

\subsection{Conditions de terminaison}

\begin{itemize}
    \item $r \leq r_s$: absorption par l'horizon $\rightarrow$ pixel noir.
    \item $r \geq 30\,r_s$: évasion vers l'infini $\rightarrow$ couleur du fond.
    \item Traversée du plan $y = 0$ dans $[r_{\text{inner}},\, r_{\text{outer}}]$
          $\rightarrow$ couleur du disque d'accrétion.
    \item $n > 500$ pas: rayon piégé près de la sphère de photons $\rightarrow$ pixel noir. 
\end{itemize}

\section{Modèle du disque d'accrétion}

\subsection{Géométrie}

Le disque est modélisé comme un anneau plan dans le plan équatorial $y = 0$,
entre le rayon de l'orbite stable la plus interne (\textbf{ISCO}) et un rayon
externe:

\begin{align}
    r_{\text{inner}} &= 3\,r_s \quad \text{(ISCO en Schwarzschild)} \\
    r_{\text{outer}} &= 12\,r_s
\end{align}

La traversée est détectée par un changement de signe de la coordonnée $y$
entre deux pas d'intégration consécutifs:

\begin{equation}
    y_{n-1} \cdot y_n < 0
    \quad \text{et} \quad
    r_{\text{inner}} < r_n < r_{\text{outer}}
\end{equation}

\subsection{Profil de température}

La température décroît avec le rayon. On paramétrise par:

\begin{equation}
    t = 1 - \frac{r - r_{\text{inner}}}{r_{\text{outer}} - r_{\text{inner}}} \in [0, 1]
\end{equation}

Le gradient de couleur va du blanc-bleu chaud ($t = 1$, près de l'ISCO)
à l'orange-rouge froid ($t = 0$, bord extérieur).

\subsection{Effet Doppler relativiste}

La matière du disque suit une orbite képlérienne de vitesse:

\begin{equation}
    v_{\text{orb}} = \sqrt{\frac{r_s}{2r}}
\end{equation}

La direction orbitale au point d'intersection $\vec{p}$ est:

\begin{equation}
    \hat{v}_{\text{orb}} = \frac{\vec{e}_y \times \hat{p}}{|\vec{e}_y \times \hat{p}|}
\end{equation}

Le \textbf{facteur Doppler} modifie la fréquence et l'intensité perçues:

\begin{equation}
    D = \frac{1}{1 - v_{\text{orb}}\cos\theta}
    \qquad \text{où} \qquad
    \cos\theta = \hat{v}_{\text{orb}} \cdot \hat{k}
\end{equation}

$\hat{k}$ étant la direction du photon au point d'intersection.
L'intensité spécifique observée scale en $D^3$ (\textbf{beaming relativiste}).

\subsection{Redshift gravitationnel}

Un photon émis à rayon $r$ et observé à l'infini subit un décalage vers le rouge:

\begin{equation}
    g = \sqrt{1 - \frac{r_s}{r}}
\end{equation}

L'intensité finale combinée est:

\begin{equation}
    \boxed{I_{\text{obs}} = D^3 \cdot g \cdot I_{\text{base}}(r)}
\end{equation}

\section{Architecture GPU}

\subsection{Pipeline de rendu}

Chaque frame est produite en deux phases successives:

\begin{enumerate}
    \item \textbf{Compute pass} \--- un thread GPU par pixel intègre une géodésique
    complète et écrit la couleur résultante dans une \emph{storage texture}.
    \item \textbf{Render pass} \--- un triangle plein écran lit la texture via un
    sampler et l'affiche sur la surface.
\end{enumerate}

\subsection{Caméra orbitale}

La caméra est paramétrisée en coordonnées sphériques $(\theta, \varphi, r)$:

\begin{equation}
    \vec{r}_{\text{cam}} = r
    \begin{pmatrix}
        \sin\theta\cos\varphi \\
        \cos\theta \\
        \sin\theta\sin\varphi
    \end{pmatrix}
\end{equation}

Sa position et sa base $(\hat{\text{right}}, \hat{\text{up}}, \hat{\text{forward}})$
sont transmises au compute shader via un \emph{uniform buffer} mis à jour chaque frame.
La direction de chaque rayon est construite depuis les coordonnées UV normalisées du pixel:

\begin{equation}
    \hat{d} = \text{normalize}\!\left(
        u \cdot \hat{\text{right}} +
        v \cdot \hat{\text{up}} +
        \frac{\hat{\text{forward}}}{\tan(\alpha/2)}
    \right)
\end{equation}

Où $\alpha$ est l'angle d'ouverture (FOV) de la caméra.

\subsection{Complexité}

Pour une résolution $W \times H$, on lance
$\lceil W/8 \rceil \times \lceil H/8 \rceil$ workgroups de $8 \times 8$ threads.
Chaque thread exécute au plus 500 itérations RK4, soit au pire
$W \cdot H \cdot 500$ évaluations de $f(u)$ par frame, entièrement parallélisées sur le GPU.

\newpage

\begin{thebibliography}{9}

\bibitem{mtw}
Misner, C.W., Thorne, K.S., Wheeler, J.A. (1973).
\textit{Gravitation}. W.H. Freeman.

\bibitem{wald}
Wald, R.M. (1984).
\textit{General relativity}. University of Chicago Press.

\bibitem{luminet}
Luminet, J.P. (1979).
Image of a spherical black hole with thin accretion disk.
\textit{Astronomy \& Astrophysics}, 75, 228--235.

\bibitem{james}
James, O. et al. (2015).
Gravitational lensing by spinning black holes in astrophysics,
and in the movie Interstellar.
\textit{Classical and Quantum Gravity}, 32(6).

\bibitem{wgpu}
wgpu \--- Safe and portable GPU abstraction in Rust.
\url{https://wgpu.rs} 

\end{thebibliography}

\end{document}
